\subsection*{Local feedback leaves useful contact-network content without an advantage over force compensation}

An independently supported elastic network tests whether the equivalence above is peculiar to a lattice transport coefficient. This auxiliary host contains fifteen existing material units connected only by actual near contact. Changing a connection's rest length changes the forces on its two endpoints. It is not a cellular automaton, has no tissue birth, and receives a fixed left support and a periodic rightward load. Its task is resource uptake under a spatially shifted supply, not acquisition of autonomous locomotion. A standard local perturbation rule adjusts rest lengths from the uptake already received at the connection's endpoints, without a supplied target direction or resource-band location.

Across three main histories, plastic rest lengths yield mean uptake 2.27404 during the first 48-step stopped-writing window and 0.75034 after the resource shift, compared with 2.09823 and 0.29758 for fixed parameters. A stronger control, however, uses fixed reference springs plus local compensating forces and matches the plastic body's complete recorded state within $8.9\times10^{-16}$. Plastic bodies also end with mean energy 50.42232, compared with 65.40015 for the fixed-parameter control. Increased uptake therefore supplies neither a net resource advantage nor additional capability beyond the force-equivalent readout. The main plastic candidate creates no new contacts after its initial network forms and breaks none; sign updates and passive creep can disconnect material despite total mass remaining fifteen (Supplementary Tables~\ref{tab:contact-stages}--\ref{tab:contact-ledger-extra}).

Three further histories compare long-term updating with the same random material perturbations but no long-term update. All three first hold windows favour updating, by 0.04866, 0.19341 and 0.06673 uptake units. After the supply shift, the differences are 0.11448, $-0.10138$ and 0.02805: only two remain positive. Six complete acquired/revised parents then supply thirteen dependent probes. Removing only the long-term rest-length offset reduces uptake in five parents and improves it in one. Preserving the old effective forces in a local compensating readout restores uptake in all six. The intervention preserves position and current perturbation initially but pays the positive change in elastic energy and a parameter-change charge, so it also changes endpoint energy; it is not a mechanically neutral memory lesion. The exact sham, independent content-use comparisons and force compensation together support a small useful contribution from ordinary locally acquired content, without establishing robust revision, learned topology or an extra material method (Supplementary Tables~\ref{tab:contact-pairs}--\ref{tab:contact-null}).


\subsection*{Contact feedback can arrive without being used effectively}

A further six-batch experiment separates access to consequences from the ability to exploit them in the same class of moving-node elastic host. Resource uptake at distal material emits a signal that travels at most one actual contact per update. In the opportunity run, the first uptake occurs at postclock 53 and the first proximal signal at 56, although the proximal material never takes up resource itself. Cutting only information transmission across a section in three actual parents preserves every recorded physical variable during 96 stopped-writing updates. The residual local signal is not erased. With writing enabled, the cut-minus-open uptake changes are tiny and of mixed sign; nearly full distal stores prevent this comparison from establishing a useful communication benefit (Supplementary Tables~\ref{tab:reach-means} and \ref{tab:reach-cut}).

Useful action consequences nevertheless exist. Three new bodies receive an actual 96-step resource-free prefix, then paid positive or negative proximal rest-length changes. Positive-minus-negative uptake is positive in all three bodies at resource position 5.4 (0.44420--0.45451) and negative at position 4.7 ($-0.42971$ to $-0.41614$). The received local signals retain the corresponding ordering. These are executed fixed-action comparisons, not counterfactual labels supplied to the learner, and they establish neither online credit assignment nor a unique causal role for one contact (Supplementary Table~\ref{tab:reach-actions}).

Providing reachable consequences does not make the tested ordinary rule reliably useful. Mean-age alignment lowers changed-environment uptake in all three histories compared with ordinary diffusion (means 0.75176 versus 1.23979). Holding actions for 96 instead of 24 updates improves that uptake in only one of three further histories; its higher terminal energy accompanies fewer paid parameter changes. Finally, three independent sparse-resource histories keep distal stores below capacity, yet diffusion yields less changed-environment uptake than no signal in all three, by 0.00542, 0.01881 and 0.00635 units. Instantaneous global feedback is not consistently better either. These unequal-information controls expose a failure of utilization under the tested rule, rather than a general absence of accessible consequences or a fair ranking of learning algorithms. Fixed material with a compensating readout exactly reproduces the first-batch physical states. Neither a new transmission-method advantage nor a broad impossibility of contact-mediated learning follows (Supplementary Tables~\ref{tab:reach-pairs}--\ref{tab:reach-jobs}).

